%% kapa: exam=1 round=21 year=2010 period=1 form=A subject=civil_law subject_ko=민법 seq=35 number=35
%% source: Q-Net 감정평가사 기출문제 (cst00309) artlSeq=5024181 file=extracted/1008/5024181/2010년 제21회 감정평가사 1차시험 기출문제/1교시.hwp
%% conversion: kordoc 4.24.51 -> tools/gen1.py (auto); review=proofread-sonnet; flags=none
\begin{question}{35}
2008년 1월 5일 甲은 친구 乙과 명의신탁약정을 한 후, 乙을 내세워 두 사람 사이에 명의신탁약정이 있다는 사실을 모르는 丙으로부터 토지를 매수하면서(매수자금은 甲이 제공) 乙명의로 소유권이전등기를 마쳤다. 나아가 甲은 乙과 통모하여 형식상 임대차계약서를 작성하여 위 토지의 임차인인 것처럼 가장하여 지금까지 그 토지를 점유ㆍ사용하여 왔다. 다음 설명으로 옳지 않은 것은? (다툼이 있으면 판례에 의함)
\begin{choices}
  \item 甲과 乙 사이의 명의신탁 약정은 무효이다.
  \item 乙은 위 부동산의 소유권을 취득한다.
  \item 甲은 乙에게 토지 그 자체가 아니라 토지의 매수자금 상당액을 부당이득으로 반환청구할 수 있다.
  \item 甲과 乙 사이의 임대차계약은 통정허위표시에 해당하여 무효이다.
  \item 甲은 乙로부터 부당이득을 반환받을 때까지 위 토지에 대하여 유치권을 행사할 수 있다.
\end{choices}
\end{question}
\answer{5}
